Wir betrachten die differenzierbare Abbildung
\maabbeledisp {\varphi} {\R^3} {\R
} {(x,y,z)} { x^2+y^4+z^6-1
} {.}
Die Menge $M$ ist die Faser von $\varphi$ über $0$. Es ist

\mathdisp {\left(D\varphi\right)_{(x,y,z)}  = (2x,4y^3,6z^5)} { . }

Diese Ableitung ist nur bei
\mathl{(x,y,z)=(0,0,0)}{} gleich $(0,0,0)$, und dies ist kein Punkt von $M$, sodass $\varphi$ in jedem Punkt von $M$ regulär ist. Daher liegt nach dem Satz über implizite Abbildungen eine zweidimensionale Mannigfaltigkeit vor.

Als Faser einer stetigen Abbildung ist $M$ eine abgeschlossene Teilmenge von $\R^3$. Ferner ist $M$ beschränkt. Für
\mathl{(x,y,z) \in M}{} ist nämlich
\mathl{\betrag {  x  } , \betrag {  y  } ,\betrag {  z  } \leq 1}{,} da andernfalls
\mathl{x^2+y^4+z^6 >1}{} wäre. Dies impliziert die Kompaktheit.

 [[Kategorie:Latexseite]]
